\begin{frame}[fragile]{附录 A：中英混排、公式与代码字体}
\begin{columns}[T,onlytextwidth]
\begin{column}{.49\textwidth}
\subhead{正文层次}
\textbf{研究对象：GPU 上的批量并行计算}\par\vspace{3mm}
Global memory、Warp shuffle、Register pressure 和 Occupancy 使用一致的术语。
\par\vspace{4mm}{\itshape Source Sans 3 italic}\par
{\bfseries Source Sans 3 bold / 思源黑体粗体}\par\vspace{4mm}
\captiontext{图注用于解释符号、条件与单位，重要信息随图表阅读。}
\end{column}
\begin{column}{.46\textwidth}
\subhead{数学与等宽字符}
\[\mathbf y=\mathbf A\mathbf x,\quad \hat a_k=\sum_{j=0}^{n-1}a_j\omega^{jk}\]
\begin{lstlisting}
size_t index = block * stride + lane;
// O0  Il1  {} [] ()  <= >= !=
\end{lstlisting}
\smalltext\texttt{JetBrains Mono: 0123456789}\par\vspace{3mm}
\notebox{中文正文使用思源黑体；西文使用 Source Sans 3；代码使用 JetBrains Mono。}
\end{column}
\end{columns}
\end{frame}

\begin{frame}{附录 B：注记色块随文字高度自适应}
\begin{columns}[T,onlytextwidth]
\begin{column}{.47\textwidth}
\subhead{简短注记}
\notebox{计时包含设备端同步。}
\vspace{5mm}\subhead{两行注记}
\notebox{此处 32 lanes 都参与 vote。Sign 的整轮 rejection 与这里的系数筛选分属不同层次。}
\vspace{5mm}\subhead{关键结论}
\insight{让注记紧邻它所解释的内容。}
\end{column}
\begin{column}{.47\textwidth}
\subhead{多行注记}
\notebox{当说明包含条件、适用范围或计量单位时，左侧色块按整个文本盒的高度绘制。色块宽度为 2 pt，文字左侧间距为 1.5 mm，适用于不同列宽和自然换行。}
\vspace{5mm}\subhead{内容与层次}
注记用于补充影响理解的条件；浅蓝结论区用于承载本页要点。
\par\vspace{3mm}\captiontext{重要条件随图表展示；文献、许可与详细证据放入随附文档或附录。}
\end{column}
\end{columns}
\end{frame}

\begin{frame}{附录 C：长标题能够自然换行，\protect\\并保持内容区域与页面边界清晰}
\subhead{阅读顺序明确的三线表}
\begin{tabularx}{\textwidth}{@{}p{2.5cm}XX@{}}
\toprule 内容类型 & 展示方式 & 常用命令 \\\midrule
概念与条件 & 小标题、正文、注记 & \texttt{\string\subhead}, \texttt{\string\notebox} \\
公式与符号 & 双栏与符号表 & \texttt{columns}, \texttt{tabularx} \\
数值与趋势 & 矢量图、图注、指标 & \texttt{\string\fig}, \texttt{\string\metric} \\
过程与实现 & Mermaid、代码片段 & \texttt{\string\includegraphics}, \texttt{lstlisting} \\\bottomrule
\end{tabularx}
\vspace{4mm}\captiontext{示例文件按主题分组；复制需要的 frame 即可复用相应版式。}
\end{frame}

\begin{frame}[fragile]{附录 D：从最小文档开始组织自己的汇报}
\begin{columns}[T,onlytextwidth]
\begin{column}{.57\textwidth}
\lstinputlisting[language=TeX]{examples/minimal-document.tex}
\end{column}
\begin{column}{.38\textwidth}
\subhead{文件入口}
\texttt{example.tex}：完整示例。\par\vspace{3mm}
\texttt{starter.tex}：简洁起点。\par\vspace{3mm}
\texttt{fonts/}：字体及原始许可。\par\vspace{3mm}
\texttt{README.md}：构建与版式说明。
\notebox{使用 XeLaTeX 编译，并保留主题、字体与所引用资源的相对目录结构。}
\end{column}
\end{columns}
\end{frame}
